% Partie : intuitionistic-logic
% Chapitre : soundness-completeness
% Section : soundness-nd

\documentclass[../../../include/open-logic-section]{subfiles}

\begin{document}

\olfileid[fr]{int}{sc}{snd}

\olsection{Correction de la déduction naturelle}

Nous allons démontrer la correction de la déduction naturelle pour la
sémantique relationnelle : si !!a{formula} est !!{derivable} à partir
d'un ensemble d'hypothèses, alors cet ensemble entraîne la
!!{formula}.

\begin{thm}[Correction]
  \ollabel{thm:soundness}
  Si $\Gamma \Proves !A$, alors $ \Gamma \Entails !A$.
\end{thm}

\begin{proof}
  Nous démontrons que si $\Gamma \Proves !A$, alors
  $\Gamma \Entails !A$. La preuve procède par récurrence sur la
  !!{derivation} de~$!A$ à partir de~$\Gamma$.

  \begin{enumerate}
  \item Si la !!{derivation} se réduit à l'hypothèse~$!A$, nous avons
    $!A \Proves !A$ et devons montrer $!A \Entails !A$. Supposons
    $\mSat{M}{!A}[w]$. Alors, trivialement,
    $\mSat{M}{!A}[w]$.

  \item La !!{derivation} se termine par $\Intro{\land}$ :
    \iftag{probAnd}{Exercice.}{Les !!{derivation}s des prémisses
      $!B$ à partir des hypothèses !!{undischarged}s~$\Gamma$ et
      de~$!C$ à partir des hypothèses !!{undischarged}s~$\Delta$
      établissent $\Gamma \Proves !B$ et $\Delta \Proves !C$.
      Par hypothèse de récurrence, $\Gamma \Entails !B$ et
      $\Delta \Entails !C$. Il faut montrer
      $\Gamma \cup \Delta \Entails !B \land !C$, car les hypothèses
      !!{undischarged}s de la !!{derivation} entière sont celles de
      $\Gamma$ réunies à celles de $\Delta$. Supposons
      $\mSat{M}{\Gamma \cup \Delta}[w]$. Alors
      $\mSat{M}{\Gamma}[w]$ aussi. Puisque $\Gamma \Entails !B$,
      nous avons $\mSat{M}{!B}[w]$ et, de même,
      $\mSat{M}{!C}[w]$. Donc $\mSat{M}{!B \land !C}[w]$.}

  \item La !!{derivation} se termine par $\Elim{\land}$ :
    \iftag{probAnd}{Exercice.}{La !!{derivation} de la prémisse
      $!B \land !C$ à partir des hypothèses !!{undischarged}s
      $\Gamma$ établit $\Gamma \Proves !B \land !C$. Par hypothèse
      de récurrence, $\Gamma \Entails !B \land !C$. Il faut montrer
      $\Gamma \Entails !B$. Supposons $\mSat{M}{\Gamma}[w]$.
      Puisque $\Gamma \Entails !B \land !C$, nous avons
      $\mSat{M}{!B \land !C}[w]$, donc $\mSat{M}{!B}[w]$.
      De même, si $\Elim\land$ conclut~$!C$, alors
      $\Gamma \Entails !C$.}

  \item La !!{derivation} se termine par $\Intro{\lor}$ :
    \iftag{probOr}{Exercice.}{Supposons que la prémisse soit $!B$
      et que les hypothèses !!{undischarged}s de la !!{derivation}
      concluant~$!B$ soient $\Gamma$. Alors $\Gamma \Proves !B$
      et, par hypothèse de récurrence, $\Gamma \Entails !B$.
      Il faut montrer $\Gamma \Entails !B \lor !C$. Supposons
      $\mSat{M}{\Gamma}[w]$. Puisque $\Gamma \Entails !B$,
      nous avons $\mSat{M}{!B}[w]$, donc aussi
      $\mSat{M}{!B \lor !C}[w]$. De même, si la prémisse est~$!C$,
      nous avons $\Gamma \Entails !C$.}

  \item La !!{derivation} se termine par~$\Elim{\lor}$ :
    \iftag{probOr}{Exercice.}{Les !!{derivation}s concluant les
      prémisses établissent $!B \lor !C$ à partir des hypothèses
      !!{undischarged}s~$\Gamma$, $!D$ à partir des hypothèses
      !!{undischarged}s $\Delta_1 \cup \{!B\}$ et $!D$ à partir des
      hypothèses !!{undischarged}s~$\Delta_2 \cup \{!C\}$. Ainsi,
      $\Gamma \Proves !B \lor !C$,
      $\Delta_1 \cup \{!B\} \Proves !D$ et
      $\Delta_2 \cup \{!C\} \Proves !D$. Par hypothèse de récurrence,
      $\Gamma \Entails !B \lor !C$,
      $\Delta_1 \cup \{!B\} \Entails !D$ et
      $\Delta_2 \cup \{!C\} \Entails !D$. Il faut démontrer
      $\Gamma \cup \Delta_1 \cup \Delta_2 \Entails !D$.

    Supposons $\mSat{M}{\Gamma \cup \Delta_1 \cup \Delta_2}[w]$.
    Alors $\mSat{M}{\Gamma}[w]$ et, puisque
    $\Gamma \Entails !B \lor !C$, nous avons
    $\mSat{M}{!B \lor !C}[w]$. Par définition de
    $\mSat{M}{}$, on a soit $\mSat{M}{!B}[w]$, soit
    $\mSat{M}{!C}[w]$. Distinguons les deux cas.
    (a)~$\mSat{M}{!B}[w]$. Alors
    $\mSat{M}{\Delta_1 \cup \{!B\}}[w]$. Puisque
    $\Delta_1 \cup \{!B\} \Entails !D$, nous obtenons
    $\mSat{M}{!D}[w]$.
    (b)~$\mSat{M}{!C}[w]$. Alors
    $\mSat{M}{\Delta_2 \cup \{!C\}}[w]$. Puisque
    $\Delta_2 \cup \{!C\} \Entails !D$, nous obtenons
    $\mSat{M}{!D}[w]$. Dans les deux cas,
    $\mSat{M}{!D}[w]$, comme voulu.}

  \item La !!{derivation} se termine par $\Intro{\lif}$ et conclut
    $!B\lif !C$. Sa prémisse est~$!C$, et la !!{derivation} de cette
    prémisse a pour hypothèses !!{undischarged}s~$\Gamma \cup
    \{!B\}$. Ainsi, $\Gamma \cup \{!B\} \Proves !C$ et, par
    hypothèse de récurrence, $\Gamma \cup \{!B\} \Entails !C$.
    Il faut montrer $\Gamma \Entails !B \lif !C$.

    Supposons $\mSat{M}{\Gamma}[w]$. Nous voulons montrer que,
    pour tout $w'$ tel que $Rww'$, si $\mSat{M}{!B}[w']$, alors
    $\mSat{M}{!C}[w']$. Supposons donc $Rww'$ et
    $\mSat{M}{!B}[w']$. Par la \olref[sem][rel]{prop:true-monotonic},
    $\mSat{M}{\Gamma}[w']$. Puisque
    $\Gamma \cup \{!B\} \Entails !C$, nous avons
    $\mSat{M}{!C}[w']$, comme voulu.

  \item La !!{derivation} se termine par $\Elim{\lif}$ et conclut
    $!C$. Ses prémisses sont $!B \lif !C$ et $!B$, obtenues par des
    !!{derivation}s à partir des hypothèses !!{undischarged}s
    $\Gamma$ et $\Delta$. Ainsi, $\Gamma \Proves !B \lif !C$ et
    $\Delta \Proves !B$. Par hypothèse de récurrence,
    $\Gamma \Entails !B \lif !C$ et $\Delta \Entails !B$.
    Il faut montrer $\Gamma \cup \Delta \Entails !C$.

    Supposons $\mSat{M}{\Gamma \cup \Delta}[w]$. Puisque
    $\mSat{M}{\Gamma}[w]$ et $\Gamma \Entails !B \lif !C$,
    nous avons $\mSat{M}{!B \lif !C}[w]$. Par définition,
    pour tout $w'$ tel que $Rww'$, si $\mSat{M}{!B}[w']$,
    alors $\mSat{M}{!C}[w']$. Comme $R$ est réflexive,
    $w$ est parmi les $w'$ tels que $Rww'$ ; ainsi,
    $\mSat{M}{!B}[w]$ implique $\mSat{M}{!C}[w]$. Puisque
    $\mSat{M}{\Delta}[w]$ et $\Delta \Entails !B$,
    nous avons $\mSat{M}{!B}[w]$, donc $\mSat{M}{!C}[w]$,
    comme voulu.

  \item La !!{derivation} se termine par $\FalseInt$ et conclut
    $!A$. Sa prémisse est $\lfalse$ et les hypothèses
    !!{undischarged}s de la !!{derivation} de cette prémisse sont
    $\Gamma$. Alors $\Gamma \Proves \lfalse$ et, par hypothèse de
    récurrence, $\Gamma \Entails \lfalse$. Il faut montrer
    $\Gamma \Entails !A$.

    Raisonnons par l'absurde. Si $\Gamma \Entails/ !A$, il existe
    un modèle~$\mModel{M}$ et un monde~$w$ tels que
    $\mSat{M}{\Gamma}[w]$ et $\mSat/{M}{!A}[w]$. Puisque
    $\Gamma \Entails \lfalse$, nous aurions
    $\mSat{M}{\lfalse}[w]$. C'est impossible, car, par définition,
    $\mSat/{M}{\lfalse}[w]$. Donc $\Gamma \Entails !A$.
  \item La !!{derivation} se termine par $\Intro\lnot$ : exercice.
  \item La !!{derivation} se termine par $\Elim\lnot$ : exercice.
  \end{enumerate}
\end{proof}

\begin{prob}
  Complétez la preuve du \olref[int][sc][snd]{thm:soundness}. Pour les
  cas de $\Intro\lnot$ et $\Elim\lnot$, utilisez la définition de
  $\mSat{M}{\lnot !A}[w]$ dans la
  \olref[int][sem][rel]{defn:true-at-w} : ne traitez pas
  $\lnot !A$ comme une abréviation de $!A \lif \lfalse$.
\end{prob}

\begin{prob}
  Montrez que les !!{formula}s suivantes ne sont pas !!{derivable}s
  en logique intuitionniste :
  \begin{enumerate}
    \item $(!A \lif !B) \lor (!B \lif !A)$
    \item $(\lnot\lnot !A \lif !A) \lif (!A \lor \lnot !A)$
    \item $(!A \lif !B \lor !C) \lif \bigl((!A \lif !B)\lor(!A \lif !C)\bigr)$
  \end{enumerate}
\end{prob}

\end{document}
